% GENERATED FILE. Edit canonical lessons, then run npm run generate:books.
% canonical-source-hash: fnv1a64:06888ce80dd5ccfa
% canonical-lessons: PA-C223-time-return, PA-C223-next-week, PA-C223-shaid, PA-C223-spoken-plan, PA-C223-payoff

\chapter{Perhaps Next Week}
\label{ch:pa-a2-plan}
\hlchaptermodality{223}

\begin{chapteropening}
\textbf{By the end of this chapter:} \emph{I can tell a short plan for tomorrow or next week, and say perhaps with shaid.}

Tell a four-line plan with a time word first and the future last, hedged with shaid.
\end{chapteropening}

\section[\texorpdfstring{\pa{ਅਗਲੇ} \pa{ਹਫ਼ਤੇ} · \pa{ਅਗਲੀ} … (agle hafte …)}{ਅਗਲੇ ਹਫ਼ਤੇ · ਅਗਲੀ … (agle hafte …)}]{Return to next week}

\label{lesson:PA-C223-time-return}



\begin{quote}
\emph{Hear:} \emph{agle hafte} Say \textbf{\pa{ਅਗਲੇ} \pa{ਹਫ਼ਤੇ}}, \textquotedblleft{}next week\textquotedblright{}. Then \emph{the next journey}: \textbf{\pa{ਅਗਲੀ} \pa{ਯਾਤਰਾ}}.
\end{quote}

\begin{grammarlens}[title={Grammar Lens}]
You can now say \emph{we will go} and \emph{I will do}. A time word at the front tells when: \textbf{\pa{ਕੱਲ੍ਹ}}, \textbf{\pa{ਅਗਲੇ} \pa{ਹਫ਼ਤੇ}}.
\end{grammarlens}

\subsection*{Your turn}
\begin{itemize}
\raggedright
  \item \emph{Hear:} \emph{next week} Say \textbf{\pa{ਅਗਲੇ} \pa{ਹਫ਼ਤੇ}}.
  \item \emph{Hear:} \emph{we will go tomorrow} Say \textbf{\pa{ਅਸੀਂ} \pa{ਕੱਲ੍ਹ} \pa{ਜਾਵਾਂਗੇ}}.
  \item \emph{Hear:} \emph{the journey will be by train} Say \textbf{\pa{ਯਾਤਰਾ} \pa{ਰੇਲ} \pa{ਨਾਲ} \pa{ਹੋਵੇਗੀ}}.
\end{itemize}

\begin{tcolorbox}[breakable,colback=teal!4,colframe=teal!35!black,title={Before you move on}]
Say \emph{next week} and \emph{tomorrow} once more.
\end{tcolorbox}
\section[\texorpdfstring{\pa{ਅਗਲੇ} \pa{ਹਫ਼ਤੇ} \pa{ਅਸੀਂ} … (agle hafte asīṁ …)}{ਅਗਲੇ ਹਫ਼ਤੇ ਅਸੀਂ … (agle hafte asīṁ …)}]{\pa{ਅਗਲੇ} \pa{ਹਫ਼ਤੇ} \pa{ਅਸੀਂ} \pa{ਬਾਜ਼ਾਰ} \pa{ਜਾਵਾਂਗੇ} — next week we will go to the market}

\label{lesson:PA-C223-next-week}



\begin{quote}
Say \emph{next week}, then \emph{we will go to the market}.
\end{quote}

\begin{grammarlens}[title={Grammar Lens: when, then who, then the future}]
A plan opens with its time, then names who, and keeps the future verb last:

\begin{quote}
\textbf{\pa{ਅਗਲੇ} \pa{ਹਫ਼ਤੇ} \pa{ਅਸੀਂ} \pa{ਬਾਜ਼ਾਰ} \pa{ਜਾਵਾਂਗੇ।}} — \emph{agle hafte asīṁ bāzār jāvāṅge.} — Next week we will go to the market.
\end{quote}

\begin{quote}
\textbf{\pa{ਅਗਲੀ} \pa{ਯਾਤਰਾ} \pa{ਰੇਲ} \pa{ਨਾਲ} \pa{ਹੋਵੇਗੀ।}} — \emph{aglī yātrā rel nāl hovegī.} — The next journey will be by train.
\end{quote}

\begin{quote}
\textbf{\pa{ਕੱਲ੍ਹ} \pa{ਸਵੇਰੇ} \pa{ਮੈਂ} \pa{ਕੰਮ} \pa{ਕਰਾਂਗਾ।}} — \emph{kallh savere maiṁ kamm karāṅgā.} — Tomorrow morning I will work.
\end{quote}
\end{grammarlens}

\subsection*{Your turn}
\begin{itemize}
\raggedright
  \item \emph{Hear:} \emph{next week Aman will go home} Say \textbf{\pa{ਅਗਲੇ} \pa{ਹਫ਼ਤੇ} \pa{ਅਮਨ} \pa{ਘਰ} \pa{ਜਾਵੇਗਾ।}}
  \item \emph{Hear:} \emph{tomorrow morning we will work} Say \textbf{\pa{ਕੱਲ੍ਹ} \pa{ਸਵੇਰੇ} \pa{ਅਸੀਂ} \pa{ਕੰਮ} \pa{ਕਰਾਂਗੇ।}}
  \item \emph{Read it:} \textbf{\pa{ਅਗਲੀ} \pa{ਯਾਤਰਾ} \pa{ਰੇਲ} \pa{ਨਾਲ} \pa{ਹੋਵੇਗੀ।}} Give its meaning.
  \item \emph{Write it:} copy \textbf{\pa{ਅਗਲੇ} \pa{ਹਫ਼ਤੇ} \pa{ਅਸੀਂ} \pa{ਬਾਜ਼ਾਰ} \pa{ਜਾਵਾਂਗੇ।}} once
\end{itemize}

\begin{tcolorbox}[breakable,colback=teal!4,colframe=teal!35!black,title={Before you move on}]
Say one plan for next week, time first.
\end{tcolorbox}
\section[\texorpdfstring{\pa{ਸ਼ਾਇਦ} (shāid)}{ਸ਼ਾਇਦ (shāid)}]{\pa{ਸ਼ਾਇਦ} (shāid) — perhaps}

\label{lesson:PA-C223-shaid}



\begin{quote}
Before the new word: say \emph{next week we will go to the market}, then \emph{Aman will come}.
\end{quote}

\subsection*{You'll want to know: \pa{ਸ਼ਾਇਦ}}
\textbf{\pa{ਸ਼ਾਇਦ}} — \emph{shāid} — \textquotedblleft{}perhaps\textquotedblright{}. It usually opens the sentence and softens a plan:

\begin{quote}
\textbf{\pa{ਸ਼ਾਇਦ} \pa{ਅਮਨ} \pa{ਕੱਲ੍ਹ} \pa{ਆਵੇਗਾ।}} — \emph{shāid Aman kallh āvegā.} — Perhaps Aman will come tomorrow.
\end{quote}

\begin{cousinweb}
\textbf{\pa{ਸ਼ਾਇਦ}} (\emph{shāid}) is a loan: it is usually traced to Persian \emph{shāyad}, \textquotedblleft{}it may be; perhaps\textquotedblright{}. CDIAL gives no Punjabi form.
\end{cousinweb}

\begin{grammarlens}[title={What you've built}]
One more word for talking about plans.
\end{grammarlens}

\subsection*{Your turn}
\begin{itemize}
\raggedright
  \item \emph{Say it:} \emph{shāid}
  \item \emph{Hear:} \emph{perhaps Aman will come tomorrow} Say \textbf{\pa{ਸ਼ਾਇਦ} \pa{ਅਮਨ} \pa{ਕੱਲ੍ਹ} \pa{ਆਵੇਗਾ।}}
  \item \emph{Recall:} say \emph{next week}, then \emph{perhaps}
  \item \emph{Write it:} copy \textbf{\pa{ਸ਼ਾਇਦ}} once from the model
\end{itemize}

\begin{tcolorbox}[breakable,colback=teal!4,colframe=teal!35!black,title={Before you move on}]
What does \textbf{\pa{ਸ਼ਾਇਦ}} mean? (\textquotedblleft{}Perhaps\textquotedblright{}.) Say it once more.
\end{tcolorbox}
\section[\texorpdfstring{\pa{ਸ਼ਾਇਦ} \pa{ਅਸੀਂ} \pa{ਅਗਲੇ} \pa{ਹਫ਼ਤੇ} … (shāid asīṁ agle …)}{ਸ਼ਾਇਦ ਅਸੀਂ ਅਗਲੇ ਹਫ਼ਤੇ … (shāid asīṁ agle …)}]{A four-line plan}

\label{lesson:PA-C223-spoken-plan}



\begin{quote}
Say \emph{perhaps Aman will come tomorrow}. Then \emph{next week we will go to the market}.
\end{quote}

\begin{grammarlens}[title={Grammar Lens: sure, and perhaps}]
Four lines, each with its time first. \textbf{\pa{ਸ਼ਾਇਦ}} marks the one that is not sure:

\begin{quote}
\textbf{\pa{ਕੱਲ੍ਹ} \pa{ਮੈਂ} \pa{ਬਾਜ਼ਾਰ} \pa{ਜਾਵਾਂਗੀ।}} — \emph{kallh maiṁ bāzār jāvāṅgī.} — Tomorrow I will go to the market.
\end{quote}

\begin{quote}
\textbf{\pa{ਅਮਨ} \pa{ਕੱਲ੍ਹ} \pa{ਕੰਮ} \pa{ਕਰੇਗਾ।}} — \emph{Aman kallh kamm karegā.} — Aman will work tomorrow.
\end{quote}

\begin{quote}
\textbf{\pa{ਸ਼ਾਇਦ} \pa{ਅਸੀਂ} \pa{ਅਗਲੇ} \pa{ਹਫ਼ਤੇ} \pa{ਜਾਵਾਂਗੇ।}} — \emph{shāid asīṁ agle hafte jāvāṅge.} — Perhaps we will go next week.
\end{quote}

\begin{quote}
\textbf{\pa{ਅਗਲੀ} \pa{ਯਾਤਰਾ} \pa{ਰੇਲ} \pa{ਨਾਲ} \pa{ਹੋਵੇਗੀ।}} — \emph{aglī yātrā rel nāl hovegī.} — The next journey will be by train.
\end{quote}
\end{grammarlens}

\subsection*{Your turn}
\begin{itemize}
\raggedright
  \item \emph{Hear:} \emph{tomorrow I will go to the market} Say it with the ending that fits you.
  \item \emph{Hear:} \emph{perhaps we will go next week} Say \textbf{\pa{ਸ਼ਾਇਦ} \pa{ਅਸੀਂ} \pa{ਅਗਲੇ} \pa{ਹਫ਼ਤੇ} \pa{ਜਾਵਾਂਗੇ।}}
  \item \emph{Read it:} the four lines Say them aloud, in order.
  \item \emph{Write it:} copy the third line once
\end{itemize}

\begin{tcolorbox}[breakable,colback=teal!4,colframe=teal!35!black,title={Before you move on}]
Say the four lines again without the model.
\end{tcolorbox}
\section[\texorpdfstring{\pa{ਸ਼ਾਇਦ} \pa{ਅਸੀਂ} \pa{ਅਗਲੇ} \pa{ਹਫ਼ਤੇ} … (shāid asīṁ agle …)}{ਸ਼ਾਇਦ ਅਸੀਂ ਅਗਲੇ ਹਫ਼ਤੇ … (shāid asīṁ agle …)}]{My plan: a four-skill payoff}

\label{lesson:PA-C223-payoff}



\begin{quote}
Say \emph{next week}, \emph{tomorrow} and \emph{perhaps}.
\end{quote}

\subsection*{Your turn: four skills}
\begin{itemize}
\raggedright
  \item \emph{Hear:} \emph{next week we will go home} Say \textbf{\pa{ਅਗਲੇ} \pa{ਹਫ਼ਤੇ} \pa{ਅਸੀਂ} \pa{ਘਰ} \pa{ਜਾਵਾਂਗੇ।}}
  \item \emph{Hear:} \emph{perhaps Aman won't come} Say \textbf{\pa{ਸ਼ਾਇਦ} \pa{ਅਮਨ} \pa{ਨਹੀਂ} \pa{ਆਵੇਗਾ।}}
  \item \emph{Read it:} \textbf{\pa{ਕੱਲ੍ਹ} \pa{ਸਵੇਰੇ} \pa{ਅਸੀਂ} \pa{ਕੰਮ} \pa{ਕਰਾਂਗੇ।}} Give its meaning: tomorrow morning we will work.
  \item \emph{Write it:} cover the line and write \textbf{\pa{ਸ਼ਾਇਦ} \pa{ਅਸੀਂ} \pa{ਅਗਲੇ} \pa{ਹਫ਼ਤੇ} \pa{ਜਾਵਾਂਗੇ।}}
\end{itemize}

\begin{tcolorbox}[breakable,colback=teal!4,colframe=teal!35!black,title={Before you move on}]
Tell your own plan in four lines: tomorrow, then next week, with one \textbf{\pa{ਸ਼ਾਇਦ}}.
\end{tcolorbox}
